% Editable, source-faithful transcription of the complete licensed author article.
% DOI 10.4171/DM/963; Documenta Mathematica29(2024),805--814.
% Original SHA256:1b2efc7c0e578c35c2580dd82186adaadeebfc3bc9c7da0b2a1d201734c705c9
% Editorial notes are separate. No PDF/image inclusion is used for the author body.
\documentclass[10pt]{article}
\usepackage[a4paper,margin=25mm]{geometry}
\usepackage{fontspec}
\setmainfont{lmroman10-regular.otf}[BoldFont=lmroman10-bold.otf,ItalicFont=lmroman10-italic.otf,BoldItalicFont=lmroman10-bolditalic.otf]
\usepackage{amsmath,amssymb,amsthm,tikz-cd}
\usepackage{hyperref}
\hypersetup{colorlinks=true,linkcolor=blue,citecolor=blue,urlcolor=blue}
\newtheorem{theorem}{Theorem}[section]
\newtheorem{lemma}[theorem]{Lemma}
\newtheorem{proposition}[theorem]{Proposition}
\theoremstyle{definition}
\newtheorem{definition}[theorem]{Definition}
\newtheorem{remark}[theorem]{Remark}
\DeclareMathOperator{\Br}{Br}
\DeclareMathOperator{\Gal}{Gal}
\DeclareMathOperator{\res}{res}
\DeclareMathOperator{\Alg}{Alg}
\DeclareMathOperator{\Ext}{Ext}
\DeclareMathOperator{\SU}{SU}
\DeclareMathOperator{\SL}{SL}
\DeclareMathOperator{\gen}{gen}
\newcommand{\A}{\mathcal A}
\newcommand{\B}{\mathcal B}
\newcommand{\C}{\mathcal C}
\newcommand{\D}{\mathcal D}
\newcommand{\Fk}{F_{(K,L,\A)}}
\newcommand{\Kl}{K_{(L,\A)}}
\newcommand{\Fs}{F_{(K,S,A)}}
\newcommand{\Ks}{K_{(S,A)}}
\newcommand{\bthree}[1]{{}_3\Br(#1)}
\newcommand{\sourcepage}[1]{\label{source-page-#1}}
\renewcommand{\labelenumi}{(\theenumi)}
\renewcommand{\qedsymbol}{$\blacksquare$}
\setlength{\parskip}{2pt}
\setlength{\parindent}{14pt}
\begin{document}
\section*{Editorial problem 3406}
\noindent\textbf{Selected problem.} Prove the exact Theorem 1.3 of the article below:
\begin{quote}
There exists a simple simply connected algebraic group $G$ over a (certain) field $T$ that is an outer form of type $A_2$ for which the genus $\gen_T(G)$ is infinite.
\end{quote}
The source defines genus in Definition 1.1 and supplies its whole construction and proof in Sections 2--3. Its constructed field is infinitely generated and has characteristic zero. The intermediate embedding results and individual algebras are supporting material, not separate problems.

\noindent\textbf{Difficulty: H2.} The construction keeps an infinite collection of degree-three division algebras distinct while making every invariant cubic field embed in each. It builds transferred generic Severi--Brauer splitting fields without losing Brauer classes, cyclicizes cubic extensions, and uses Pfister function fields to conjugate unitary involutions. A transfinite induction and repeated finite-data descent then control all existing and newly appearing maximal invariant fields. Coordinating ramification, Brauer restriction, unitary embeddings and maximal tori puts the authored proof beyond ordinary PhD qualifying algebra.

\noindent\textbf{Source and license.} Sergey V. Tikhonov, \textit{Outer forms of type $A_2$ with infinite genus}, Documenta Mathematica \textbf{29} (2024), 805--814; DOI \href{https://doi.org/10.4171/DM/963}{10.4171/DM/963}. Copyright 2024 Deutsche Mathematiker-Vereinigung; published by EMS Press. The individual publisher page and original PDF cover grant \href{https://creativecommons.org/licenses/by/4.0/}{Creative Commons Attribution 4.0 (CC BY 4.0)}.

\noindent\textbf{Transcription origin.} The following complete article is an editable source-faithful transcription of the licensed human-authored publisher PDF, not native author TeX obtained from the publisher. The author's English, order, hypotheses, formulas, proof material, remarks and references are preserved. All three mathematical diagrams are editable TeX. Formatting is adapted; no replacement proof or new mathematical argument is added. Source-fidelity observations are kept separately in \path{../sources/3406/EDITORIAL-NOTES.md}.

\noindent\textbf{Original source.} \url{https://ems.press/journals/dm/articles/14297947}

\noindent\textbf{Pinned original PDF SHA256.}\par
\noindent\texttt{1b2efc7c0e578c35c2580dd82186adaadeeb}\par
\noindent\texttt{fc3bc9c7da0b2a1d201734c705c9}

\clearpage
\setcounter{page}{805}
\sourcepage{1}
\noindent\begin{minipage}[t]{.47\textwidth}\small
Doc. Math. 29 (2024), 805--814\\ DOI 10.4171/DM/963
\end{minipage}\hfill\begin{minipage}[t]{.52\textwidth}\small\raggedleft
\copyright\ 2024 Deutsche Mathematiker-Vereinigung\\Published by EMS Press\\
This work is licensed under a \href{https://creativecommons.org/licenses/by/4.0/}{CC BY 4.0} license
\end{minipage}
\begin{center}\vspace{12pt}
{\Large\bfseries Outer forms of type $A_2$ with infinite genus}\par\vspace{10pt}
{\large Sergey V. Tikhonov}\par\vspace{10pt}
\end{center}
\noindent\textbf{Abstract.} Let $G$ be an absolutely almost simple algebraic group over a field $K$. The genus $\gen_K(G)$ of $G$ is the set of $K$-isomorphism classes of $K$-forms $G'$ of $G$ that have the same $K$-isomorphism classes of maximal $K$-tori as $G$. We construct an example of outer forms of type $A_2$ with infinite genus.
\section{Introduction}
Let $K$ be a field and $K^{\mathrm{sep}}$ its separable closure. Two absolutely almost simple algebraic $K$-groups $G_1$ and $G_2$ are said to have the same $K$-isomorphism classes of maximal $K$-tori if every maximal $K$-torus of $G_1$ is $K$-isomorphic to some maximal $K$-torus of $G_2$, and vice versa. An algebraic $K$-group $G'$ is called a $K$-form of an algebraic $K$-group $G$ if $G$ and $G'$ are isomorphic over $K^{\mathrm{sep}}$.
\begin{definition}[{\cite[Def. 6.1]{r2}}]
Let $G$ be an absolutely almost simple algebraic group over a field $K$. The genus $\gen_K(G)$ of $G$ is the set of $K$-isomorphism classes of $K$-forms $G'$ of $G$ that have the same $K$-isomorphism classes of maximal $K$-tori as $G$.
\end{definition}
The genus is trivial in some special cases and it is conjectured to be finite whenever the field $K$ is finitely generated of ``good'' characteristic (see details in \cite[\S8]{r6}).

In a similar way one can define the genus of a division algebra.
\begin{definition}
The genus $\gen(\D)$ of a finite-dimensional central division algebra $\D$ over a field $K$ is defined as the set of classes $[\D']\in\Br(K)$, where $\D'$ is a central division $K$-algebra having the same maximal subfields as $\D$.
\end{definition}
If $\D$ is a finite-dimensional central division $K$-algebra, then it is well known that any maximal $K$-torus of the corresponding algebraic group $G=\SL_{1,\D}$ is of the form $R_{E/K}(\mathbb G_m)\cap G$ (where $R_{E/K}(\mathbb G_m)$ is the Weil restriction of the $1$-dimensional split torus $\mathbb G_m$) for some maximal separable subfield $E$ of $\D$. Thus the results on genus of division algebras from \cite{r4,r10} rephrased in the language of algebraic groups say that for any prime $p$, there exist fields (with infinite transcendence degree over the prime subfield) over which there are inner forms of type $A_{p-1}$ with infinite genus. An example of groups

\medskip\noindent\small\textit{Mathematics Subject Classification 2020:} 16K20 (primary); 20G15 (secondary).\\
\textit{Keywords:} division algebra, algebraic group, genus, maximal subfield of a division algebra, algebra with involution.\normalsize
\newpage
\sourcepage{2}
of type $G_2$ with infinite genus is obtained in \cite[Rem. 3.6 (b)]{r1}. In the present paper, we construct such an example for outer forms of type $A_2$.

Let $F/K$ be a quadratic separable field extension and $\sigma$ the non-trivial $K$-automorphism of $F$. An involution on an $F$-algebra $R$ is called an $F/K$-involution if its restriction to $F$ is $\sigma$. An isomorphism of $F$-algebras with involution
\[
f:(R,\tau)\longrightarrow(R',\tau')
\]
is an $F$-algebra isomorphism $f:R\to R'$ such that $\tau'\circ f=f\circ\tau$. Let also $E/F$ be a field extension such that $E$ has an automorphism $\psi$ of order $2$ such that $\psi|_F=\sigma$ (i.e., $E$ has an $F/K$-involution). Then $\tau_E$ denotes the involution on $R\otimes_F E$ defined by the formula $\tau_E(r\otimes e):=\tau(r)\otimes\psi(e)$, where $r\in R$, $e\in E$. In particular, if $L/F$ is a field extension linearly disjoint to $E$ over $F$ with an automorphism $\phi$ of order two extending $\sigma$, then $\phi_E$ is the automorphism of order two of the field $EL=L\otimes_F E$ which extends $\phi$.

Let $\A$ be a central division $F$-algebra of degree $n$ with an $F/K$-involution $\tau$. Let $L/F$ be a separable field extension of degree $n$, and let $\phi:L\to L$ be an automorphism of order two such that $\phi|_F=\sigma$. An embedding of algebras with involution $(L,\phi)\hookrightarrow(\A,\tau)$ is by definition an injective $F$-homomorphism $f:L\to\A$ such that $\tau\circ f=f\circ\phi$. It is known that embeddings of maximal tori into the special unitary group $\SU(\A,\tau)$ can be described in terms of embeddings of fields with involution into the central simple algebra with involution $(\A,\tau)$ \cite[Prop. 2.3]{r5}.

In this paper, we construct a field $E$ and a subfield $T\subset E$ such that $[E:T]=2$ and there is an infinite set of (pairwise non-isomorphic) division $E$-algebras of degree $3$ with $E/T$-involution such that a field extension $L/E$ of degree $3$ can be embedded as an algebra with involution into one algebra of this set if and only if it can be embedded as an algebra with involution into all other algebras from this set. Passing to the corresponding special unitary groups, we obtain an example of outer forms of type $A_2$ with infinite genus. The main result of the paper is the following theorem.
\begin{theorem}\label{thm-main}
There exists a simple simply connected algebraic group $G$ over a (certain) field $T$ that is an outer form of type $A_2$ for which the genus $\gen_T(G)$ is infinite.
\end{theorem}
Note that the field $T$ constructed in the proof of this theorem is infinitely generated.

Below we use the following notation: $\Alg_3(F/K)$ is the set of isomorphism classes of central division $F$-algebras of degree $3$ with $F/K$ involution; $\Ext_3(F/K)$ is the set of isomorphism classes of field extensions of $F$ of degree $3$ with $F/K$-involution. The $3$-torsion of the Brauer group $\Br(F)$ is denoted by $\bthree F$. For a field extension $E/F$ and a central simple $F$-algebra $\A$, $\A_E$ denotes the tensor product $\A\otimes_F E$ and $\res_{E/F}:\Br(F)\to\Br(E)$ denotes the restriction homomorphism. The restriction of $\res_{E/F}$ to the subgroup $\bthree F$ will also be denoted by $\res_{E/F}$. For a central simple $F$-algebra $\A$, $\A^{\mathrm{op}}$ denotes the opposite algebra and $\A^m$ denotes $\A\otimes_F\cdots\otimes_F\A$ ($m$ times). For a quadratic form $q$ over $K$ and a field extension $E/K$, $q_E$ denotes the quadratic form obtained by extension of scalars from $K$ to $E$. Recall that a field extension $E/F$ is called regular if $E/F$ is separable and $F$ is algebraically closed in $E$.
\newpage
\sourcepage{3}
\section{Preliminary results}
We start with the following.
\begin{lemma}\label{lem-transfer}
Let $n$ be a positive integer, $F$ a field of characteristic not dividing $2n$, $F/K$ a quadratic field extension, $\sigma$ the non-trivial $K$-automorphism of $F$, $\A$ a central simple $F$-algebra of degree $n$, and $L/F$ a cyclic field extension of degree $n$. Then there exists a regular field extension $M/F$ and a subfield $T\subset M$ such that $[M:T]=2$ and
\begin{enumerate}
\item $M=TF$ and the non-trivial $T$-automorphism of $M$ extends $\sigma$;
\item the composite $ML$ splits $\A_M$;
\item the homomorphism $\res_{M/F}:\Br(F)\to\Br(M)$ is injective.
\end{enumerate}
\end{lemma}
\begin{proof}
Let $F(x)$ be a purely transcendental extension of $F$ of transcendence degree $1$. Let also $\phi$ be a generator of the Galois group $\Gal(L(x)/F(x))$ and
\[
\C:=\A^{\mathrm{op}}_{F(x)}\otimes_{F(x)}(L(x)/F(x),\phi,x),
\]
where $(L(x)/F(x),\phi,x)$ is a cyclic $F(x)$-algebra of degree $n$. Let also $E$ be the function field of the Severi--Brauer variety of $\C$. Note that the kernel of the restriction homomorphism $\res_{E/F(x)}:\Br(F(x))\to\Br(E)$ is generated by $[\C]$ (see, e.g., \cite[Cor. 13.16]{r8}).

Let $\B$ be a central simple $F$-algebra of exponent bigger than $1$. Assume that $\B$ is split by $E$, then $[\B_{F(x)}]=[\C^i]$ for some $1\leq i\leq n$. If $i<n$, then the $F(x)$-algebra $\C^i$ ramifies at the discrete valuation (trivial on $F$) of $F(x)$ defined by the polynomial $x$, but $\B_{F(x)}$ is unramified at this valuation, hence $[\B_{F(x)}]\neq[\C^i]$. Since the exponent of $\B_{F(x)}$ is bigger than $1$, then
\[
[\B_{F(x)}]\neq[\C^n]=[F(x)].
\]
Thus $\B_{F(x)}$ is not split by $E$, i.e., the homomorphism $\res_{E/F}:\Br(F)\to\Br(E)$ is injective.

Since $E$ splits $\C$, then
\[
[\A_E]=[(L(x)/F(x),\phi,x)_E]=[(EL/E,\phi',x)],
\]
where $\phi'$ is the generator of $\Gal(EL/E)$. Thus $EL$ splits $\A_E$.

Note that $E/F$ is a regular extension of $F(x)$. For the following construction of the transfer of a regular field extension, we refer to \cite[p. 220]{r7}.

Let $\sigma$ also denotes the $K(x)$-automorphism of $F(x)$ extending the automorphism $\sigma$ of $F$. The automorphism $\sigma$ of $F(x)$ can be extended to an isomorphism (which we also denote by $\sigma$) of $E$ and another regular extension of $F(x)$ denoted by $E_\sigma$. Thus the following diagram commutes:
\begin{equation}\label{eq-transfer}
\begin{tikzcd}[column sep=large,row sep=large]
F\arrow[r,hook]\arrow[d,"\sigma"'] & F(x)\arrow[r,hook]\arrow[d,"\sigma"'] & E\arrow[d,"\sigma"]\\
F\arrow[r,hook] & F(x)\arrow[r,hook] & E_\sigma
\end{tikzcd}\tag{2.1}
\end{equation}
\newpage
\sourcepage{4}
Let $M=EE_\sigma$ be the free composite over $F$ of $E$ and $E_\sigma$. This free composite is $F$-isomorphic to the function field of the Severi--Brauer variety of the $E_\sigma(y)$-algebra
\[
\A^{\mathrm{op}}_{E_\sigma(y)}\otimes_{E_\sigma(y)}(E_\sigma L(y)/E_\sigma(y),\psi',y),
\]
where $y$ is transcendental over $E_\sigma$ (we replace $x$ by $y$ since the composite is free) and $\psi'$ is the generator of the Galois group $\Gal(E_\sigma L(y)/E_\sigma(y))$. The field $M$ is a regular extension of $F$. The isomorphisms
\[
\sigma:E\longrightarrow E_\sigma\quad\text{and}\quad\sigma^{-1}:E_\sigma\longrightarrow E
\]
have a unique extension to an automorphism $\bar\sigma$ of $M$ of order two. Let $T:=T_{F/K}(E)$ be the transfer of $E$ with respect to the ground field descent $F\supset K$, i.e., $T$ is the subfield of $M$ of elements fixed under the action of $\bar\sigma$. Note that the composite $TF$ coincides with $M$, $[M:T]=2$, and $\bar\sigma$ extends $\sigma$.

The algebra $\A_M$ is split by $ML$ since $\A_E$ is split by $EL$.

Finally, the diagram \eqref{eq-transfer} induces the following commutative diagram for the corresponding Brauer groups:
\[
\begin{tikzcd}[column sep=large,row sep=large]
\Br(F)\arrow[r,"\res"]\arrow[d,"\cong"'] & \Br(F(x))\arrow[r,"\res"]\arrow[d,"\cong"'] & \Br(E)\arrow[d,"\cong"] & {}\\
\Br(F)\arrow[r,"\res"] & \Br(F(x))\arrow[r,"\res"] & \Br(E_\sigma)\arrow[r,"\res"] & \Br(M).
\end{tikzcd}
\]
The injectivity of $\res_{E/F}$ implies the injectivity of $\res_{E_\sigma/F}$. Moreover, $\res_{M/E_\sigma}$ is injective by the same arguments as for $\res_{E/F}$, we just replace the ground field $F$ by $E_\sigma$. Hence the homomorphism $\res_{M/F}$ is also injective.
\end{proof}
We also need the following.
\begin{lemma}\label{lem-cyclic}
Let $F$ be a field of characteristic $\neq2,3$; $F/K$ a quadratic field extension, $\sigma$ the non-trivial $K$-automorphism of $F$, and $L/F$ a field extension of degree $3$ with an automorphism $\phi$ of order two such that $\phi|_F=\sigma$. Then there exists a field extension $F(L)/F$ and a subfield $K(L)\subset F(L)$ such that $[F(L):K(L)]=2$ and
\begin{enumerate}
\item $F(L)=K(L)F$ and the non-trivial $K(L)$-automorphism of $F(L)$, denoted by $\sigma_{F(L)}$, extends $\sigma$;
\item $[F(L):F]\leq2$;
\item the composite $F(L)L$ is a cyclic extension of $F(L)$ of degree $3$;
\item the homomorphism $\res_{F(L)/F}:\bthree F\to\bthree{F(L)}$ is injective.
\end{enumerate}
\end{lemma}
\begin{proof}
If the extension $L/F$ is cyclic, then one can take $F(L):=F$, $K(L):=K$, and $\sigma_{F(L)}:=\sigma$.

Assume that the extension $L/F$ is not cyclic. Let $N$ be the normal closure of the extension $L^\phi/K$, where $L^\phi\subset L$ is the subfield of elements fixed by $\phi$. Then $F\not\subset N$, and
\newpage
\sourcepage{5}
$NF$ is the normal closure of the extension $L/F$. Thus we have the following diagram of field extensions:
\[
\begin{tikzcd}[column sep=large,row sep=large]
N\arrow[dr,dash,"2"'] & {} & L\arrow[dl,dash,"2"]\arrow[dr,dash,"3"] & {}\\
{} & L^\phi\arrow[dr,dash,"3"'] & {} & F\arrow[dl,dash,"2"]\\
{} & {} & K & {}
\end{tikzcd}
\]
Let $H$ be the Sylow $3$-subgroup of the Galois group $\Gal(N/K)$. Then $N^H$, the fixed field of $H$, is an extension of $K$ of degree $2$ and $N/N^H$ is a cyclic extension of degree $3$. Hence $NF$ is a cyclic extension of $N^HF$ of degree $3$ and $[N^HF:F]=2$. Let $K(L):=N^H$ and $F(L):=K(L)F$. Since $F\not\subset N$, then $[F(L):K(L)]=2$ and the field $F(L)$ has a $K(L)$-automorphism of order two extending $\sigma$. Note that $F(L)L=NF$, hence $F(L)L/F(L)$ is a cyclic extension of degree $3$. Finally, since $[F(L):F]=2$, then the homomorphism $\res_{F(L)/F}:\bthree F\to\bthree{F(L)}$ is injective.
\end{proof}
\begin{remark}
In the notations of Lemma~\ref{lem-cyclic}, for any field extension $L'$ of $F$ of degree $3$, $F(L)$ and $L'$ are linearly disjoint over $F$. Moreover, if $L'$ has an automorphism $\phi'$ of order $2$ extending $\sigma$, then the composite $F(L)L'$ has the automorphism $\phi'_{F(L)}$ which extends the automorphisms $\phi'$ and $\sigma_{F(L)}$.
\end{remark}
\begin{lemma}\label{lem-embedding}
Let $F$ be a field of characteristic $\neq2,3$; $F/K$ a quadratic field extension, $\sigma$ the non-trivial $K$-automorphism of $F$, $\A$ a central simple $F$-algebra of degree $3$ with an $F/K$-involution $\tau$, and $L/F$ a field extension of degree $3$ with an automorphism $\phi$ of order two such that $\phi|_F=\sigma$. Then there exists a field extension $\Fk/F$ and a subfield $\Kl\subset\Fk$ such that $[\Fk:\Kl]=2$ and
\begin{enumerate}
\item $\Fk=\Kl F$ and the non-trivial $\Kl$-automorphism of $\Fk$, denoted by $\sigma_{\Fk}$, extends $\sigma$;
\item the homomorphism $\res_{\Fk/F}:\bthree F\to\bthree{\Fk}$ is injective;
\item for any field extension $L'$ of $F$ of degree $3$, $\Fk$ and $L'$ are linearly disjoint over $F$;
\item there is an embedding $(\Fk L,\phi_{\Fk})\hookrightarrow(\A_{\Fk},\tau_{\Fk})$ of algebras with involution.
\end{enumerate}
\end{lemma}
\begin{proof}
Let $F(L)$, $K(L)$ and $\sigma_{F(L)}$ be as in Lemma~\ref{lem-cyclic}. Let also $M$ and $T$ be fields obtained by applying Lemma~\ref{lem-transfer} for the quadratic field extension $F(L)/K(L)$, the $F(L)$-algebra $\A_{F(L)}$, the cyclic field extension $F(L)L/F(L)$ of degree $3$. Then by Lemmas~\ref{lem-transfer} and~\ref{lem-cyclic}, the homomorphism $\res_{M/F}:\bthree F\to\bthree M$ is injective; for any field extension $L'$ of $F$ of degree $3$, $M$ and $L'$ are linearly disjoint over $F$ and the composite $ML$ splits $\A_M$. Thus there is an $M$-embedding $\varepsilon:ML\hookrightarrow\A_M$ of $M$-algebras.
\newpage
\sourcepage{6}
Note that $\A_M$ has the $M/T$-involution $\tau_M$ which extends $\tau$ and $ML$ has the automorphism $\phi_M$ of order two extending $\phi$. By \cite[Proposition 3.1]{r5}, there exists an $M/T$-involution $\delta$ on $\A_M$ such that $\varepsilon:(ML,\phi_M)\hookrightarrow(\A_M,\delta)$ is an embedding of algebras with involution.

Let $\pi(\delta)$ and $\pi(\tau_M)$ be the $3$-fold Pfister forms of involutions $\delta$ and $\tau_M$ respectively (see \cite[\S19.B]{r3}). Let $T(\pi(\delta))$ and $T(\pi(\tau_M))$ be the function fields of $\pi(\delta)$ and $\pi(\tau_M)$ respectively. Then the quadratic forms $\pi(\delta)_{T(\pi(\delta))}$ and $\pi(\tau_M)_{T(\pi(\tau_M))}$ are isotropic and hence hyperbolic since they are Pfister forms.

Let $\Kl$ be the free composite over $T$ of the fields $T(\pi(\delta))$ and $T(\pi(\tau_M))$. Let also $\Fk:=\Kl F$. Since $F\not\subset\Kl$, then $[\Fk:\Kl]=2$ and $\Fk$ has a $\Kl$-automorphism $\sigma_{\Fk}$ of order $2$ extending $\sigma$. Note that the algebraic closure of $F$ in $\Fk$ is $F(L)$ and $[F(L):F]$ is either $1$ or $2$. Therefore, $\Fk$ and $L'$ are linearly disjoint over $F$ for any field extension $L'$ of $F$ of degree $3$.

The extensions $T(\pi(\delta))/T$ and $T(\pi(\tau_M))/T$ are composition of a purely transcendental extension with a quadratic extension, hence the homomorphism
\[
\res_{\Fk/M}:\bthree M\longrightarrow\bthree{\Fk}
\]
is injective. Hence $\res_{\Fk/F}:\bthree F\to\bthree{\Fk}$ is also injective.

The quadratic forms $\pi(\delta)_{\Kl}$ and $\pi(\tau_M)_{\Kl}$ are hyperbolic. Then by \cite[Th. 19.6]{r3}, the involutions $\delta_{\Fk}$ and $\tau_{\Fk}$ on $\A_{\Fk}$ are conjugate. This means that there is an isomorphism $\xi:(\A_{\Fk},\delta_{\Fk})\to(\A_{\Fk},\tau_{\Fk})$ of algebras with involution.

Moreover, the embedding $\varepsilon:(ML,\phi_M)\hookrightarrow(\A_M,\delta)$ of algebras with involution induces an embedding $(\Fk L,\phi_{\Fk})\hookrightarrow(\A_{\Fk},\delta_{\Fk})$ of algebras with involution. Indeed, $\Fk L=ML\otimes_M\Fk$. Let
\[
\varepsilon_{\Fk}:ML\otimes_M\Fk\longrightarrow\A_{\Fk}
\]
be an $\Fk$-embedding defined by the formula $\varepsilon_{\Fk}(m\otimes a):=\varepsilon(m)\otimes a$, where $m\in ML$, $a\in\Fk$. Then
\begin{align*}
\varepsilon_{\Fk}\bigl(\phi_{\Fk}(m\otimes a)\bigr)
 &=\varepsilon_{\Fk}\bigl(\phi_M(m)\otimes\sigma_{\Fk}(a)\bigr)\\
 &=\varepsilon\bigl(\phi_M(m)\bigr)\otimes\sigma_{\Fk}(a)\\
 &=\delta\bigl(\varepsilon(m)\bigr)\otimes\sigma_{\Fk}(a)\\
 &=\delta_{\Fk}\bigl(\varepsilon(m)\otimes a\bigr)\\
 &=\delta_{\Fk}\bigl(\varepsilon_{\Fk}(m\otimes a)\bigr).
\end{align*}
Thus $\varepsilon_{\Fk}$ is an embedding of algebras with involutions. Then $\xi\circ\varepsilon_{\Fk}$ is an embedding
\[
(\Fk L,\phi_{\Fk})\hookrightarrow(\A_{\Fk},\tau_{\Fk})
\]
of algebras with involution.
\end{proof}
The following construction of the field $\Fs$ is an adaptation of the construction from \cite{r10} for algebras with involutions. We give the details below for the reader's convenience.
\newpage
\sourcepage{7}
\begin{proposition}\label{prop-transfinite}
Let $F$ be a field of characteristic $\neq2,3$; $F/K$ a quadratic field extension, $\sigma$ the non-trivial $K$-automorphism of $F$, $A\subset\Alg_3(F/K)$ and $S\subset\Ext_3(F/K)$. Then there exists a field extension $\Fs/F$ and a subfield $\Ks\subset\Fs$ such that $[\Fs:\Ks]=2$ and
\begin{enumerate}
\item $\Fs=\Ks F$ and the non-trivial $\Ks$-automorphism, denoted by $\sigma_{\Fs}$, of $\Fs$ extends $\sigma$;
\item the homomorphism $\res_{\Fs/F}:\bthree F\to\bthree{\Fs}$ is injective;
\item for any field extension $L'$ of $F$ of degree $3$, $\Fk$ and $L'$ are linearly disjoint over $F$;
\item for any $\A\in A$ with an $F/K$-involution $\tau$ and $L\in S$ with a $K$-automorphism $\phi$ of order $2$ extending $\sigma$, there is an embedding
\[
(\Fs L,\phi_{\Fs})\hookrightarrow(\A_{\Fs},\tau_{\Fs})
\]
of algebras with involution.
\end{enumerate}
\end{proposition}
\begin{proof}
Let $\mathcal P:=\{(L,\D)\mid L\in S\text{ and }\D\in A\}$ be the set of pairs. Let also $<$ be a well-ordering on $\mathcal P$ and let $t_0=(L_0,\D_0)$ denote its least element. Set $E_{t_0}:=F_{(K,L_0,\D_0)}$ and $T_0:=K_{(L_0,\D_0)}$, where the fields $F_{(K,L_0,\D_0)}$ and $K_{(L_0,\D_0)}$ are constructed in Lemma~\ref{lem-embedding}. For $t=(L,\D)\in\mathcal P$, set
\begin{gather*}
E^{<t}:=\bigcup_{t'<t}E_{t'},\qquad T^{<t}:=\bigcup_{t'<t}T_{t'},\\
T_t:=T^{<t}_{(E^{<t}L,\D_{E^{<t}})},\qquad
E_t:=E^{<t}_{(T^{<t},E^{<t}L,\D_{E^{<t}})},
\end{gather*}
where the fields $E_t$ and $T_t$ are obtained by applying Lemma~\ref{lem-embedding} to the quadratic field extension $E^{<t}/T^{<t}$, the field extension $E^{<t}L/E^{<t}$ of degree $3$, the automorphism $\phi_{E^{<t}}$ of $E^{<t}L$ extending the automorphism $\phi$ of $L$ and the $E^{<t}$-algebra $\D_{E^{<t}}$. We also define $\Fs:=\bigcup_{t\in\mathcal P}E_t$ and $\Ks:=\bigcup_{t\in\mathcal P}T_t$.

By Lemma~\ref{lem-embedding}, $E_t=T_tF$ and $[E_t:T_t]=2$ for any $t\in\mathcal P$. Then $\Fs=\Ks F$, $[\Fs:\Ks]=2$ and the non-trivial $\Ks$-automorphism of $\Fs$ extends $\sigma$.

By Lemma~\ref{lem-embedding} and transfinite induction, the homomorphism
\[
\res_{\Fs/F}:\bthree F\longrightarrow\bthree{\Fs}
\]
is injective and for any field extension $L'$ of $F$ of degree $3$, $\Fs$ and $L'$ are linearly disjoint over $F$.

Finally, let $\A\in A$ with an $F/K$-involution $\tau$, $L\in S$ with an automorphism $\phi$ of order $2$ extending $\sigma$ and $t=(L,\A)$. By Lemma~\ref{lem-embedding}, there is an embedding $(E_tL,\phi_{E_t})\hookrightarrow(\A_{E_t},\tau_{E_t})$ of algebras with involution. Moreover, as in the proof of Lemma~\ref{lem-embedding}, this embedding induces the embedding
\[
(\Fs L,\phi_{\Fs})\hookrightarrow(\A_{\Fs},\tau_{\Fs})
\]
of algebras with involution.
\end{proof}
\newpage
\sourcepage{8}
\begin{theorem}\label{thm-exhaust}
Let $F$ be a field of characteristic $\neq2,3$, $F/K$ a quadratic field extension, $\sigma$ the non-trivial $K$-automorphism of $F$, $A\subset\Alg_3(F/K)$. Then there exists a field extension $F_A/F$ and a subfield $K_A\subset F_A$ such that $[F_A:K_A]=2$ and
\begin{enumerate}
\item $F_A=K_AF$ and the non-trivial $K_A$-automorphism, denoted by $\sigma_{F_A}$, of $F_A$ extends $\sigma$;
\item the homomorphism $\res_{F_A/F}:\bthree F\to\bthree{F_A}$ is injective;
\item for any central simple $F$-algebra $\B$ of degree $3$ with an $F/K$-involution $\theta$, the algebra $\B_{F_A}$ has an $F_A/K_A$-involution $\sigma_{F_A}$ extending $\theta$;
\item if $L\in\Ext_3(F_A/K_A)$ with a $K_A$-automorphism $\phi$ of order $2$ extending $\sigma_{F_A}$, then there is an embedding
\[
(L,\phi)\hookrightarrow(\A_{F_A},\tau_{F_A})
\]
of algebras with involution for any $\A\in A$ with an $F/K$-involution $\tau$.
\end{enumerate}
\end{theorem}
\begin{proof}
Let $F_0:=F$ and $K_0:=K$. We recursively define $F_i$ and $K_i$, $i\in\mathbb Z_{>0}$, to be the fields
\[
F_{i-1\,(K_{i-1},\Ext_3(F_{i-1}/K_{i-1}),\res_{F_{i-1}/F}(A))}
\]
and
\[
K_{i-1\,(\Ext_3(F_{i-1}/K_{i-1}),\res_{F_{i-1}/F}(A))}
\]
constructed by applying Proposition~\ref{prop-transfinite} to the quadratic field extension $F_{i-1}/K_{i-1}$, the set
\[
\res_{F_{i-1}/F}(A)\subset\Alg_3(F_{i-1}/K_{i-1})
\]
and the set $\Ext_3(F_{i-1}/K_{i-1})$.

Let $F_A:=\bigcup_{i\geq0}F_i$ and $K_A:=\bigcup_{i\geq0}K_i$. Hence $F_A=K_AF$ and the non-trivial $K_A$-automorphism $\sigma_{F_A}$ of $F_A$ extends $\sigma$. Therefore, for any central simple $F$-algebra $\B$ of degree $3$ with an $F/K$-involution $\theta$, the $F_A/K_A$-involution $\sigma_{F_A}$ extends $\theta$.

By induction and Proposition~\ref{prop-transfinite}, $\res_{F_A/F}:\bthree F\to\bthree{F_A}$ is injective.

Assume that $\A\in A$ with an $F/K$-involution $\tau$ and $L\in\Ext_3(F_A/K_A)$ with an automorphism $\phi$ of order two extending $\sigma_{F_A}$. Then there exists $i\geq0$ and a field extension $L'$ of $F_i$ of degree $3$ such that $L=F_AL'$ and $\phi_i:=\phi|_{L'}$ is a $K_i$-automorphism of order two. This means that $L'\in\Ext_3(F_i/K_i)$. By Proposition~\ref{prop-transfinite}, there is an embedding
\[
(F_{i+1}L',\phi_{i\,F_{i+1}})\hookrightarrow(\A_{F_{i+1}},\tau_{F_{i+1}})
\]
of algebras with involution. As in the proof of Lemma~\ref{lem-embedding}, this embedding can be extended to an embedding $(L,\phi)\hookrightarrow(\A_{F_A},\tau_{F_A})$ of algebras with involution.
\end{proof}
\section{Construction}
As a corollary to Theorem~\ref{thm-exhaust}, we obtain the following.
\begin{theorem}\label{thm-construction}
There exists a field $E$ and a subfield $T\subset E$ with $[E:T]=2$ such that there is an infinite set $B$ of pairwise non-isomorphic division $E$-algebras of degree $3$ with $E/T$-involution and such that for any field extension $L/E$ of degree $3$ with an automorphism $\phi$ of order $2$ extending the non-trivial $T$-automorphism of $E$, there is an embedding $(L,\phi)\hookrightarrow(\A,\tau)$ of algebras with involution for any $\A\in B$ with an $E/T$-involution $\tau$.
\end{theorem}
\newpage
\sourcepage{9}
\begin{proof}
Let $\xi_3$ be a primitive $3$th root of unity, $K=\mathbb Q(\xi_3)(x,y,z)$, the purely transcendental extension of the field $\mathbb Q(\xi_3)$ and $F=K(\sqrt2)$. Then for $i>0$, the symbol $F$-algebras $\A_i=(\frac{x+\sqrt2y^i}{x-\sqrt2y^i},z)_3$ of degree $3$ are pairwise non-isomorphic. Indeed, $\A_i$ ramifies at the discrete valuation (trivial on $\mathbb Q(\xi_3)(\sqrt2)(y,z)$) of $F=\mathbb Q(\xi_3)(\sqrt2)(y,z)(x)$ defined by the polynomial $x+\sqrt2y^i\in\mathbb Q(\xi_3)(\sqrt2)(y,z)[x]$, but if $i\neq j$, then $\A_j$ is unramified at this valuation.

By the projection formula for the corestriction \cite[Th. 3.2]{r9}, we have that the corestriction to $K$ of the $F$-algebra $\A_i$ is similar to the split $K$-algebra
\[
\left(N_{F/K}\left(\frac{x+\sqrt2y^i}{x-\sqrt2y^i}\right),z\right)_3=(1,z)_3.
\]
Then by \cite[Th. 3.1 (2)]{r3}, the algebra $\A_i$ has an $F/K$-involution.

Now we apply Theorem~\ref{thm-exhaust} for the infinite set $A\subset\Alg_3(F/K)$ consisting of algebras $\A_i$, $i>0$, and set $E:=F_A$, $T:=K_A$, $B:=\res_{F_A/F}(A)$.
\end{proof}
\begin{remark}
The referee asked the following interesting question. The construction presented in the proof of the previous theorem yields algebras that are not isomorphic even as algebras without involution. One may wonder if one can construct (infinite) families of non-equivalent involutions supported on the same division algebra with isomorphic maximal invariant subfields.
\end{remark}
Now we are in a position to prove the main Theorem~\ref{thm-main}.
\begin{proof}
We use notation from Theorem~\ref{thm-construction}. Let $\A$ be an algebra with $E/T$-involution $\tau$ from the set $B$. Then the special unitary group $G=\SU(\A,\tau)$ is a simple simply connected outer form of type $A_2$ over $T$.

If $\A_1$ are $\A_2$ are different algebras with $E/T$-involutions $\tau_1$ and $\tau_2$ respectively from the set $B$, then the algebraic groups $\SU(\A_1,\tau_1)$ and $\SU(\A_2,\tau_2)$ are not isomorphic by \cite[Th. 26.9]{r3}. Moreover, by \cite[Prop. 2.3]{r5} the groups $\SU(\A_1,\tau_1)$ and $\SU(\A_2,\tau_2)$ have the same $T$-isomorphism classes of maximal $T$-tori.

Thus, the genus $\gen_T(G)$ is infinite.
\end{proof}
\noindent\textbf{Acknowledgments.} The author thanks the referee for helpful suggestions.
\begin{thebibliography}{10}
\bibitem{r1} C. Beli, P. Gille, and T.-Y. Lee, \textit{Examples of algebraic groups of type $G_2$ having the same maximal tori}. Tr. Mat. Inst. Steklova \textbf{292} (2016), 16--25; translation in Proc. Steklov Inst. Math. \textbf{292} (2016), 10--19 Zbl 1356.20028 MR 3628451
\bibitem{r2} V. I. Chernousov, A. S. Rapinchuk, and I. A. Rapinchuk, \textit{Division algebras with the same maximal subfields}. Uspekhi Mat. Nauk \textbf{70} (2015), no. 1(421), 89--122; translation in Russian Math. Surveys \textbf{70} (2015), no. 1, 83--112 Zbl 1325.12006 MR 3353117
\newpage
\sourcepage{10}
\bibitem{r3} M.-A. Knus, A. Merkurjev, M. Rost, and J.-P. Tignol, \textit{The book of involutions}. Amer. Math. Soc. Colloq. Publ. 44, American Mathematical Society, Providence, RI, 1998 Zbl 0955.16001 MR 1632779
\bibitem{r4} J. S. Meyer, \textit{Division algebras with infinite genus}. Bull. Lond. Math. Soc. \textbf{46} (2014), no. 3, 463--468 Zbl 1301.16021 MR 3210701
\bibitem{r5} G. Prasad and A. S. Rapinchuk, \textit{Local-global principles for embedding of fields with involution into simple algebras with involution}. Comment. Math. Helv. \textbf{85} (2010), no. 3, 583--645 Zbl 1223.11047 MR 2653693
\bibitem{r6} A. S. Rapinchuk and I. A. Rapinchuk, \textit{Linear algebraic groups with good reduction}. Res. Math. Sci. \textbf{7} (2020), no. 3, article no. 28 Zbl 1466.20032 MR 4149442
\bibitem{r7} P. Roquette, \textit{Isomorphisms of generic splitting fields of simple algebras}. J. Reine Angew. Math. \textbf{214(215)} (1964), 207--226 Zbl 0219.16023 MR 0166215
\bibitem{r8} D. J. Saltman, \textit{Lectures on division algebras}. CBMS Reg. Conf. Ser. Math. 94, Published by American Mathematical Society, Providence, RI; on behalf of Conference Board of the Mathematical Sciences, Washington, DC, 1999 Zbl 0934.16013 MR 1692654
\bibitem{r9} J.-P. Tignol, \textit{On the corestriction of central simple algebras}. Math. Z. \textbf{194} (1987), no. 2, 267--274 Zbl 0595.16012 MR 0876236
\bibitem{r10} S. V. Tikhonov, \textit{Division algebras of prime degree with infinite genus}. Tr. Mat. Inst. Steklova \textbf{292} (2016), 264--267; translation in Proc. Steklov Inst. Math. \textbf{292} (2016), 256--259 Zbl 1356.16014 MR 3628465
\end{thebibliography}
\bigskip\noindent Communicated by Ulf Rehmann

\noindent Received 1 September 2023; revised 26 November 2023.
\bigskip

\noindent\textbf{Sergey V. Tikhonov}\\
Department of Higher Algebra, Belarusian State University, Nezavisimosti Avenue 4,\\
220030 Minsk, Belarus; \href{mailto:tikhonovsv@bsu.by}{tikhonovsv@bsu.by}
\end{document}
